% ═══════════════════════════════════════════════════════════════
% Section 4: TARA Threat Taxonomy
% ═══════════════════════════════════════════════════════════════

\section{TARA: Threat Taxonomy}
\label{sec:tara}

The Therapeutic Applications \& Risk Assessment (TARA) registry catalogues BCI
attack techniques with a structure inspired by MITRE
ATT\&CK~\cite{mitre2024attack}, independently developed to cover neural,
cognitive, and physical safety domains. Each
technique is simultaneously an attack vector, an ethical risk, and---where
applicable---a therapeutic application.

\paragraph{Why not ATT\&CK directly?}
MITRE ATT\&CK is designed for enterprise IT and mobile environments where
adversary objectives center on data exfiltration, lateral movement, and
persistence. BCI threats differ in three fundamental ways that make direct
adoption inadequate: (1)~the target is biological tissue, not an information
system---attacks can cause seizures, tissue necrosis, or cognitive coercion,
none of which map to ATT\&CK's impact categories; (2)~the same technique is
built on the same mechanism as a therapy (e.g., the stimulation used in deep
brain stimulation for Parkinson's disease can also be delivered without authorisation), requiring a
dual-use classification that ATT\&CK has no mechanism to express; and (3)~the
attack lifecycle must span from radio-frequency transmission through silicon
processing to neural tissue, crossing the bio-digital boundary that ATT\&CK's
purely digital model does not address. TARA adopts ATT\&CK's proven structural
methodology---techniques organized into tactics and domains---while building an
independent taxonomy purpose-built for the BCI threat landscape.

\subsection{Taxonomy Structure}

The TARA registry organizes \SnapshotTechniques{} techniques into \SnapshotTactics{} tactics across \SnapshotDomains{} operational
domains. Technique identifiers follow the format \texttt{QIF-TXXX}, where the
numeric suffix provides sequential ordering.

\paragraph{Domains.}
The eight domains span the full attack lifecycle:

\begin{enumerate}
  \item \textbf{Neural (N)} --- Direct interaction with neural tissue
        (scan, injection, manipulation)
  \item \textbf{BCI System (B)} --- System-level intrusion and evasion
  \item \textbf{Protocol (P)} --- Protocol disruption and communication attacks
  \item \textbf{Data (D)} --- Data harvesting and exfiltration
  \item \textbf{Cognitive (C)} --- Cognitive exploitation and imprinting
  \item \textbf{Countermeasure (M)} --- Surveillance and monitoring
  \item \textbf{Evasion (E)} --- Defense evasion and anti-detection
  \item \textbf{Sensor (S)} --- Consumer device side-channel attacks
\end{enumerate}

\paragraph{Tactics.}
The 15 tactics follow a lifecycle structure analogous to ATT\&CK, adapted for BCI operations:
Neural Scan, BCI Intrusion, Neural Injection, Cognitive Imprinting, BCI Evasion,
Data Harvesting, Neural Manipulation, Evasion/Rootkit Deployment, Monitoring/Surveillance,
Protocol Disruption, Cognitive Exploitation, Signal Replay, Signal Harvesting,
Signal Fingerprinting, and Signal Chaining.

\subsection{Evidence Classification}

Each technique carries an evidence status. Table~\ref{tab:status} gives the snapshot
counts used in this paper beside the live counts.

\begin{table}[H]
\centering
\caption{TARA evidence status: the v\SnapshotRegistryVersion{} snapshot analysed here and
the live registry (v\LiveRegistryVersion{}, read 2026-10-17).}
\label{tab:status}
\small
\begin{tabularx}{\textwidth}{@{}l X r r@{}}
\toprule
\textbf{Status} & \textbf{Definition (live registry)} & \textbf{v\SnapshotRegistryVersion{}} & \textbf{Live} \\
\midrule
Confirmed    & Retired 2026-10-16 (see text)                                          & 19 & --- \\
Demonstrated & A published source executes the technique against a real system      & 35 & 48 \\
Emerging     & Shown against an analogous system or in part                         & 22 & 27 \\
Theoretical  & The mechanism is established but no attack using it has been shown   & 31 & 95 \\
Plausible    & The mechanism is contested or indirect; the attack is reasoned        & 1  & 2 \\
Speculative  & Neither mechanism nor attack is established                          & 1  & 1 \\
\midrule
             & \textbf{Total}                                                       & \SnapshotTechniques{} & \LiveTechniques{} \\
\bottomrule
\end{tabularx}
\end{table}

\paragraph{CONFIRMED retired.} The status \emph{confirmed} was retired on 2026-10-16
after an audit found that it spanned unrelated kinds of evidence; for example, it had been
given to signal injection (QIF-T0001) on the strength of a position paper and a P300
probing study, neither of which demonstrates signal injection. The 27 entries that carried
it in the live registry were reclassified as theoretical (21), demonstrated (5) and
plausible (1), each keeping its prior value and a reason. The 19 \emph{confirmed} entries of
the snapshot should be read with this in mind. \textbf{No status value means that an attack
has been carried out on a person with an implanted or worn device.}

\paragraph{Evidence rule.} Every live entry records whether its sources are peer-reviewed,
preprint or absent (158, 7 and 8 entries). \textbf{A technique with no peer-reviewed source
may not hold a status above theoretical and does not carry a NISS score.} Papers that test a
defence are held in a separate field, so they cannot satisfy the citation requirement for
an attack. When the rule was applied, 34 entries broke it; 24 were given a verified
peer-reviewed source and 10 lost their scores. ``Peer-reviewed'' currently means that no
preprint identifier was found, which is weaker than verified.

\paragraph{Hazards held separately.} Some conditions harm a device user with no attacker,
such as a manufacturer withdrawing support for an implant. They do not belong in a catalogue
of adversary actions, so 23 entries are held in \texttt{datalake/qif-hazards.json}. One,
QIF-H0001 (loss of manufacturer or sponsor support), rests on two peer-reviewed
sources~\cite{higgins2026posttrial,zhao2026consentinflux} and is theoretical. The other 22 are unsourced pointers
to catalogue records that describe a therapy or physical mechanism; they are not warnings
about those therapies and assert nothing about incidence in people.

\subsection{Severity Distribution}

Each technique also carries a filed severity rating. In v1.5 this table was labelled as
CVSS v4.0 base ratings; it is in fact the registry's \texttt{severity} field, assigned by
hand and not computed from the CVSS vectors. In the live registry it agrees with the NISS
severity for 30 of \LiveNissScored{} scored techniques, so it should not be read as either
a CVSS or a NISS result.

\begin{table}[H]
\centering
\caption{TARA filed severity distribution (v\SnapshotRegistryVersion{} snapshot; hand-assigned).}
\label{tab:severity}
\small
\begin{tabular}{@{}l r r@{}}
\toprule
\textbf{Severity} & \textbf{Count} & \textbf{Percentage} \\
\midrule
Critical & 29 & 26.6\% \\
High     & 60 & 55.0\% \\
Medium   & 17 & 15.6\% \\
Low      &  3 &  2.8\% \\
\bottomrule
\end{tabular}
\end{table}

The weighting toward high and critical reflects the author's judgement at filing, not a
computed result.

\subsection{Category Breakdown}

The \SnapshotTechniques{} techniques distribute across eight operational categories:

\begin{table}[H]
\centering
\caption{TARA techniques by operational category.}
\label{tab:categories}
\small
\begin{tabular}{@{}l l r@{}}
\toprule
\textbf{ID} & \textbf{Category} & \textbf{Count} \\
\midrule
SE & Signal Eavesdropping    & 20 \\
CI & Cognitive Interference  & 18 \\
EX & Data Exfiltration       & 19 \\
DM & Data Manipulation       & 15 \\
SI & Signal Injection        & 13 \\
PE & Privilege Escalation    &  9 \\
DS & Denial of Service       &  8 \\
PS & Physical Safety         &  7 \\
\bottomrule
\end{tabular}
\end{table}

\subsection{Dual-Use Mapping}

TARA assesses each technique for therapeutic analogs. The same physical mechanisms that
enable attacks---electromagnetic stimulation, signal decoding, neuromodulation---are
the mechanisms underlying established therapies such as deep brain stimulation
(DBS)~\cite{lozano2019dbs}, transcranial magnetic stimulation
(TMS)~\cite{hallett2007tms}, and neurofeedback.

\begin{table}[H]
\centering
\caption{Dual-use classification of \SnapshotTechniques{} TARA techniques.}
\label{tab:dualuse}
\small
\begin{tabular}{@{}l l r r@{}}
\toprule
\textbf{Classification} & \textbf{Definition} & \textbf{Count} & \textbf{\%} \\
\midrule
Confirmed    & Published clinical use exists      & 53 & 48.6\% \\
Probable     & Under active clinical investigation & 16 & 14.7\% \\
Possible     & Theoretical therapeutic mapping     &  9 &  8.3\% \\
Silicon Only & No tissue analog; purely digital    & 25 & 22.9\% \\
Not Mapped   & Dual-use assessment pending         &  6 &  5.5\% \\
\bottomrule
\end{tabular}
\end{table}

Of the \SnapshotTechniques{} techniques, 78 (71.6\%) have a confirmed, probable or
possible therapeutic analog, and 25 are silicon-only. The finding is that many attack
mechanisms are shared with therapies, not that every attack has a therapeutic counterpart:
the mechanism is shared, and the boundary between therapy and attack is consent, dosage and
oversight. Defences cannot simply remove a capability that treatment depends on.

\subsection{Representative Techniques}

Table~\ref{tab:representative} shows five registry entries spanning categories and dual-use
classes. In v1.5 this table listed descriptions and identifiers that did not match the
registry entries carrying those identifiers (for example, QIF-T0015 was described as a P300
side channel but is directed-energy thermal damage); it now quotes the registry.

\begin{table}[H]
\centering
\caption{Five representative TARA entries as stored in the registry (status: snapshot
$\to$ live).}
\label{tab:representative}
\small
\begin{tabularx}{\textwidth}{@{}l X l l X@{}}
\toprule
\textbf{ID} & \textbf{Name} & \textbf{Status} & \textbf{Dual-use} & \textbf{Therapeutic analog} \\
\midrule
T0001 & Signal injection & Confirmed $\to$ Theoretical & Confirmed & tDCS/tACS neuromodulation \\
T0035 & P300 interrogation & Demonstrated & Confirmed & P300 BCI for locked-in patients \\
T0058 & Calibration poisoning & Demonstrated & Probable & Adaptive BCI calibration \\
T0072 & Transducer inversion (speaker-to-microphone) & Confirmed $\to$ Demonstrated & Silicon only & --- \\
T0090 & WiFi CSI passive body sensing & Demonstrated & Confirmed & Contactless vital-sign monitoring \\
\bottomrule
\end{tabularx}
\end{table}

